% !TEX root = ../main.tex
\chapter{Forward and Backward in a Single Graph}
\label{ch:fwd_bwd_compile}

One of the six locked design decisions
(\chref{ch:compile_design}~Section~\ref{ssec:compile_design:decision_fwd_bwd_single})
was to capture forward and backward into a single \MPSGraph,
not just the forward. This chapter explains why, how the
unification works, and the trade-offs it imposes.

\section{Why Forward + Backward Together}
\label{sec:fwd_bwd_compile:why}

A training step consists of three parts: forward, loss
computation, backward. The intermediate tensors produced by the
forward (the activations) are consumed by the backward (to
compute parameter gradients). If forward and backward are
compiled separately:

\begin{itemize}
  \item Each activation must materialise (the forward compiles
    cannot know which activations the backward will need).
  \item The optimiser between the two cannot fuse across
    boundaries.
  \item The handoff incurs memory writes and synchronisation.
\end{itemize}

A unified forward+backward compile lets the optimiser see
everything: which activations are read by the backward
(materialise only those), which are not (fuse away), where
allocations can be eliminated, where memory can be reused. The
empirical win is substantial: roughly half of the compile
speedup over eager is attributable to this fusion.

\section{The MPSGraph Backward}
\label{sec:fwd_bwd_compile:mpsgraph_backward}

\MPSGraph\ added a
\code{gradientForPrimaryTensor:withTensors:name:} method in
macOS~13. Given a scalar output tensor and a list of input
tensors, it constructs the gradient tensors within the same
graph:

\begin{lstlisting}[style=lucidcpp,language=C++]
NSDictionary<MPSGraphTensor*, MPSGraphTensor*>* grads =
    [graph gradientForPrimaryTensor:loss_tensor
                        withTensors:@[param1, param2, ...]
                               name:nil];
// grads[param1] is the gradient with respect to param1
\end{lstlisting}

The internal \MPSGraph\ machinery performs reverse-mode AD on
its own DAG, exactly analogous to \Lucid's eager autograd
engine (\chref{ch:autograd_engine}). The resulting gradient
operations are added to the same graph.

The compile then optimises forward and backward together: dead
intermediates are pruned, common subexpressions are eliminated
across the boundary, layout changes are unified.

\section{How Lucid Uses It}
\label{sec:fwd_bwd_compile:how_lucid_uses}

The compile path's unified forward+backward construction:

\begin{enumerate}
  \item Trace the forward as usual (\chref{ch:tracer}). The
    trace ends at the loss tensor.
  \item Identify the parameters (the leaf tensors with
    \code{requires\_grad=True}).
  \item Emit the forward into \MPSGraph\ via the MpsBuilder
    (\chref{ch:mps_builder}).
  \item Call \code{gradientForPrimaryTensor} on the loss tensor
    with the parameter list. \MPSGraph\ adds the backward
    operations.
  \item Compile the combined graph.
\end{enumerate}

The compiled executable has the same inputs as the forward (the
training-step inputs) and additional outputs for the parameter
gradients.

\subsection{The optimiser step}
\label{ssec:fwd_bwd_compile:optimiser_step}

Should the optimiser step also be compiled into the same graph?
\Lucid's compile path does not currently do this. The reasoning:

\begin{itemize}
  \item Optimisers (Adam, SGD with momentum) maintain per-
    parameter state buffers that must persist across calls. The
    compiled graph would need to handle these as mutable
    arguments, which complicates the cache key.
  \item The optimiser step is fast ($\sim 4$ ms) compared to
    forward+backward ($\sim 100$ ms). The compilation win on the
    optimiser step alone would be small.
\end{itemize}

The eager optimiser step runs immediately after the compiled
forward+backward returns the gradients. The cost is a
synchronisation between the two, which is amortised.

\section{Saved Tensors in the Compiled Graph}
\label{sec:fwd_bwd_compile:saved_tensors}

In eager autograd, the forward pass saves intermediate tensors
(the activations) that the backward will read. In the compiled
graph, the same dependence exists but is expressed as a DAG
edge: an activation that the backward reads is alive from the
forward that produced it to the backward that consumes it.

\MPSGraph's compiler analyses this dependence and allocates the
memory accordingly: activations that the backward does not read
are eliminated; activations that are read once are allocated
into the smallest live-set at any time; recomputable activations
(in principle) could be recomputed rather than stored, though
\MPSGraph\ does not perform automatic checkpointing.

\subsection{Memory profile}
\label{ssec:fwd_bwd_compile:memory_profile}

The compiled graph's peak memory is typically 20--30\% lower than
the eager equivalent, because the compiler eliminates
intermediate activations that the eager autograd would
materialise unnecessarily.

For training the user can layer gradient checkpointing
(\chref{ch:gradient_checkpointing}) on top: each checkpointed
region is its own compile boundary. The combination gives the
best of both worlds.

\section{Numerical Equivalence}
\label{sec:fwd_bwd_compile:numerical_equiv}

The compiled forward+backward must produce numerically the same
result as the eager forward+backward. Equivalence is enforced by
parity tests at three levels:

\begin{itemize}
  \item Per-op: each emitter's parity test
    (\chref{ch:op_emitters}).
  \item Per-layer: each standard layer (Linear, Conv,
    LayerNorm, Attention) has a parity test that compiles a
    one-layer model and checks against eager.
  \item Per-model: each model-zoo entry has a parity test that
    compiles the full forward+backward and checks loss and
    parameter gradients.
\end{itemize}

The tolerance is FP32: $10^{-5}$ relative, $10^{-4}$ absolute.
The differences are due to the compiler's reordering of
operations (floating-point non-associativity).

\section{Edge Cases}
\label{sec:fwd_bwd_compile:edge_cases}

\subsection{Multiple losses}
\label{ssec:fwd_bwd_compile:multi_loss}

A model that has multiple loss terms (multi-task training)
typically sums them into a single scalar before backward. The
compile path supports this naturally: the sum is a primitive,
the unified backward propagates gradients to all parameters that
contributed.

If the user wants per-loss gradients (rare), the compile path
falls back to eager.

\subsection{Conditionally computed losses}
\label{ssec:fwd_bwd_compile:conditional_loss}

A loss that is computed only for some samples (e.g., auxiliary
loss with a per-sample mask) is naturally expressed as
masked-mean. The mask is a constant per-call, so the compile
path handles it correctly.

\subsection{Gradient accumulation}
\label{ssec:fwd_bwd_compile:grad_accumulation}

The standard gradient-accumulation pattern (run several forward+
backward steps without optimiser updates, then step) interacts
with the compile path: each forward+backward call is the
compiled graph; the eager-side gradient accumulation accumulates
the returned gradients into the parameter's \code{.grad}.

\subsection{Skipping the backward}
\label{ssec:fwd_bwd_compile:skip_backward}

If the user runs the compiled module without subsequent
\code{loss.backward()} (i.e., inference), the gradients are
computed by the compiled graph anyway and then discarded. This
is wasteful.

The mitigation: compile the model in inference mode (no
\code{requires\_grad} on the parameters), which produces a
forward-only graph.

\section{The Backward IR}
\label{sec:fwd_bwd_compile:backward_ir}

\MPSGraph's backward produces operations that are visible in the
graph dump. \code{dump\_executable} shows them alongside the
forward operations:

\begin{verbatim}
# Forward
%0 = placeholder (input)
%1 = layer_norm(%0, gamma, beta) -> %1
%2 = matmul(%1, weight)          -> %2
%3 = softmax(%2)                 -> %3
%4 = sum(%3 * target)            -> %4 (loss)

# Backward (generated by gradientForPrimaryTensor)
%5 = grad_softmax(%4) -> %5
%6 = grad_matmul(%5, weight) -> grad_input
%7 = grad_matmul(%5, %1) -> grad_weight
%8 = grad_layer_norm(%6, ...) -> grad_input, grad_gamma, grad_beta
\end{verbatim}

The structure mirrors what eager autograd would produce; the
compile-time optimisation collapses redundant work that eager
would not have caught.

\section{Implementation Notes}
\label{sec:fwd_bwd_compile:impl_notes}

The forward+backward unification is implemented in the
MpsBuilder's \code{build()} method. After emitting the forward
graph, the builder collects the parameter tensors that the trace
identified as having \code{requires\_grad=True}, calls
\code{gradientForPrimaryTensor:withTensors:}, and adds the
returned tensor dict as additional graph outputs.

The CompiledExecutable's \code{run()} returns both the forward
outputs (the loss) and the backward outputs (the parameter
gradients), in the order specified at compile time.

\section{Summary and Forward References}
\label{sec:fwd_bwd_compile:summary}

The compile path captures forward + backward into a single
\MPSGraph, using \MPSGraph's
\code{gradientForPrimaryTensor:withTensors:} method to construct
the backward operations alongside the forward. The unified
compilation lets \MPSGraph's optimiser see the full training-step
graph, fusing across the forward/backward boundary and
eliminating intermediates that the eager autograd would have
materialised.

The next chapter, \chref{ch:compiled_module}, treats the user-
facing \code{CompiledModule} API. \chref{ch:compile_hardening}
treats the edge cases that the 3.5 release stabilised.

\begin{lucidnote}
For the user: forward+backward unification is automatic.
\code{compile(model)} produces a CompiledModule that, when
called, returns the loss tensor with a connected backward graph.
\code{loss.backward()} flows through the compiled gradient
operations.
\end{lucidnote}
